\chapter{Cohomology and Serre duality}

Before approaching this chapter, it's useful to remind all of the notations that
will be used, introduced in Chapter~\ref{chp:difforms}.

\begin{itemize}[label = --]
    \item \(\Omega^{0}{(X)}\) the set of \(C^{\infty}\) functions on
\(X\)
    \item \(\Omega^{1}{(X)}\) the set of \(C^{\infty}\) 1-forms on \(X\) 
    \item \(\Omega^{2}{(X)}\) the set of \(C^{\infty}\) 2-forms on \(X\) 
    \item \(\Omega^{{(1,0)}} = \{\omega \in \Omega^{1}{(X)} : \omega = P{(z)}\,dz\} \) 
    \item \(\Omega^{{(0,1)}} = \{\omega \in \Omega^{1}{(X)} : \omega = Q{(z)}\,d\bar{z}\} \) 
\end{itemize}

Finally, the following diagram holds:
\[
\begin{tikzcd}
	& {\Omega^0(X)} & \\
	{\Omega^{(1,0)}(X)} && {\Omega^{(0,1)}(X)} \\
	{\{0\}} & {\Omega^2(X)} & {\{0\}}
	\arrow["\partial"', from=1-2, to=2-1]
	\arrow["{\bar{\partial}}", from=1-2, to=2-3]
	\arrow["\partial"', from=2-1, to=3-1]
	\arrow["{\bar{\partial}}", from=2-1, to=3-2]
	\arrow["\partial"', from=2-3, to=3-2]
	\arrow["{\bar{\partial}}", from=2-3, to=3-3]
\end{tikzcd}
\]

\section{De Rham Cohomology}
Let \(X\) be a connected compact Riemann Surface. Then The following sequence is
a complex
\[
  \{0\} \hookrightarrow \Omega^{0}{(X)} \overset{d_{1}}{\to} \Omega^{1}{(X)}
  \overset{d_{2}}{\to} \Omega^2{(X)} \to \{0\} \quad \text{(as \(\mathbb{C}\)-vector
  spaces)}
\]
And it leads to the following cohomology spaces
\begin{itemize}[label = --]
    \item \(H^{0}_{DR} = \{f \in  \Omega^{0}{(X)} : df = 0\} \cong \mathbb{C}\) 
    \item \(H^{1}_{DR} = \{\omega \in \Omega^{1}{(X)} : d\omega = 0\} /
        \{df : f \in \Omega^{0}{(X)}\} \) 
    \item \(H^{2}_{DR} = \{\eta \in \Omega^2{(X)}\} / \{d\omega : \omega \in
        \Omega^{1}{(X)}\}  \) 
\end{itemize}

\begin{theorem}{De Rham for Surfaces}
    It holds that 
    \[
      H^{0}_{DR} \cong H^2_{DR} \cong \mathbb{C} \quad \text{ and } \quad
      \mathbb{H}^{1}_{DR} \cong \mathbb{C}^{2g}
    \]
    where \(g\) is the genus of \(X\) 
\end{theorem}

\section{Dolbeault Cohomology}

The cohomology spaces are
\begin{itemize}[label = --]
    \item \(H^{{(1,0)}}{(X)} = \{\omega \in \Omega^{{(1,0)}}{(X)} :
        \bar{\partial}\omega = 0\} \) which is exactly the set of holomorphic
        1-forms
    \item \(H^{{(0,1)}}{(X)} = \{\omega \in \Omega^{{(0,1)}}{(X)}\} /
        \{\partial f : f \in \Omega^{0}{(X)}\}\) 
\end{itemize}

And there are some observations to make.

\begin{enumerate}[label = \arabic*.]
    \item \(H^{{(1,0)}}{(X)} \hookrightarrow H^{1}_{DR}{(X)} \). Indeed, if \(\omega\) is
        a holomorphic 1-form and if \(\omega = df\), then in local coordinates
        \[
          \omega = P{(z)} \, dz = \frac{\partial f}{\partial z} dz +
          \underbrace{\frac{\partial f}{\partial \bar{z}}}_{=0}  d\bar{z}
          \implies f \text{ is holomorphic } \implies f \text{ is constant}
        \]
        hence \(\omega = df = 0\) 
    \item Similarly, setting \(\overline{H^{{(1,0)}}}{(X)} = \{\omega \in
        \Omega^{{(0,1)}} {(X)} : \partial \omega = 0\} \) the set of
        anti-holomorphic 1-forms, then we also have
    \(\overline{H^{{(1,0)}}} {(X)} \hookrightarrow H^{1}_{DR}{(X)} \) which
        is shown using the same logic.
\end{enumerate}

Indeed, there is a stronger relationship

\begin{theorem}{}
    The following equivalences hold:
    \begin{align}
      H^{1}_{DR} {(X)} = H^{{(1,0)}}{(X)} \oplus \overline{H^{{(1,0)}}}{(X)}\\
    \dim H^{{(1,0)}}{(X)} = \dim
      \overline{H^{{(1,0)}}}{(X)} = g
    \end{align}    
\end{theorem}

\begin{lemma}{}
    Assume \(\eta \in \Omega^2{(X)}\) such that \(\int_X \eta = 0\). Then there
    exists \(f \in C^{\infty}{(X)}\) such that \(\Delta f = \eta\) 
\end{lemma}
The proof is be articulated in a \emph{soft proof} with techniques used in
Hilbert spaces (Lax-Milgram) and an argument about local ellipticity.

Another possibility is to use PDEs.

\begin{proof}[End of proof]
    Fix \(\omega \in \Omega^{1}{(X)}\) such that \(d\omega = 0\). We look for
    \(\omega = df + h + g\), where \(h\) is holomorphic and \(g\)
    anti-holomorphic. Say \(\omega = \omega^{{(1,0)}} + \omega^{{(0,1)}}\).
    Since \(d\omega = 0\), then \(\bar{\partial}\omega^{{(1,0)}} = -\partial
    \omega^{{(0,1)}}\). Let \(f \in C^{\infty}{(X)}\) and write \[\omega - df =
        \underbrace{{\left(  \omega^{{(1,0)}} - \partial f \right)}}_{=:h}  +
        \underbrace{{\left( \omega^{{(0,1)}} -
    \bar{\partial} f\right)}}_{=:g}  \] 
    and we want to have \(h\) holomorphic and \(g\) anti-holomorphic. However,
    \(\bar{\partial}h = 0 \iff \bar{\partial} \omega^{{(1,0)}} -
    \bar{\partial}\partial f = 0 \iff \partial \omega^{{(0,1)}} - \partial
    \bar{\partial} f = 0\), hence all is left is to solve
    \[
      \bar{\partial} \omega^{{(1,0)}} = \bar{\partial} \partial f
    \]
    But 
    \[
      \int_X \bar{\partial} \omega^{{(1,0)}} = \int_{X} d \omega^{{(1,0)}} = 0
      \text{ by Stokes}
    \]
    Hence \(\exists f \) such that \(\Delta f = \bar{\partial}
    \omega^{{(1,0)}}\) by solving Poisson's equation
\end{proof}


\begin{proof}{}
    The proof starts by constructing the isomorphism \(\omega \mapsto
    \bar{\omega}\) which maps
    \(\omega \in H^{{(1,0)}}\) with \(\omega = P{(z)} \, dz\), to 
    \(\bar{\omega} = \bar{P}{(z)} \, d\bar{z} \in \overline{H^{{(1,0)}}}{(X)}\).
    It's easy to show that this is an anti-linear isomorphism, reaching the
    second equation as a consequence.

    % TODO ehh??

    In conclusion, we have \(\dim H^{{(1,0)}}{(X)} = g\) 
\end{proof}

\begin{proposition}{}
    On \(H^{{(1,0)}}{(X)}\), 
    \[
      \Span{\omega_{1}, \omega_{2}} := -\frac{1}{2i} \int_X \omega_{1} \wedge
      \overline{\omega_{2}}
    \]
    is an inner product
\end{proposition}
\begin{proof}{}
    To show that is \emph{positive definite}, note that locally \(-\frac{1}{2i}
    \omega \wedge \bar{\omega} = |P|^2 dx \wedge dy\), hence \(\Span{\omega,
    \omega} = 0\) implies \(\omega = 0\). Everything else is obvious.
\end{proof}

\subsection{Serre's Duality}
\begin{theorem}{}\label{thm:serre}
    \[
      \left[ \beta \right] = 0 \text{ in } H^{{(0,1)}}{(X)} \iff \forall \alpha
      \in H^{{(1,0)}}{(X)}, \; \int_X \alpha \wedge \beta = 0
    \]
\end{theorem}
\begin{corollary}{}
    \[
      H^{{(0,1)}}{(X)} \cong {\left( H^{{(1,0)}}{(X)} \right)}^{*} \quad \text{
      and } \quad \dim H^{{(0,1)}}{(X)} = g
    \]
\end{corollary}
\begin{proof}[Proof of corollary]
    Note that \begin{align*}
        \zeta : H^{{(0,1)}} &\longrightarrow {\left( H^{{(1,0)}} \right)} ^{*} \\
        [\beta] &\longmapsto \zeta ([\beta]) = \left(L_\beta : \alpha \mapsto \int_X
        \alpha \wedge \beta \right)
    \end{align*}
    By the theorem, \(\zeta\) is injective. By Riesz on \({(H^{{(1,0)}},
    \Span{\cdot , \cdot } )}\), for any \(\varphi  \in {\left( H^{{(1,0)}}{(X)}
    \right)}^{*} \), there exist \(\omega_{0} \in H^{{(1,0)}}\) such that
    \[
      \varphi {(\omega)} = \Span{\omega, \omega_{0}} = -\frac{1}{2i} \int_X
      \omega \wedge \overline{\omega_{0}} = \int_X \omega \wedge \left[
      -\frac{1}{2i} \overline{\omega_{0}} \right] 
    \]
\end{proof}
\begin{proof}[Proof of Theorem~\ref{thm:serre}]
    \begin{itemize}
        \item[\(\implies \)] Obvious by Stokes
        \item[\(\impliedby \)] Let \([\beta] \in H^{{(0,1)}}{(X)}\) such that
            \(\forall \alpha \in H^{{(0,1)}}{(X)}\), \(\int_X \alpha \wedge
            \beta = 0\). Note that \(d \beta = \partial \beta\), hence \(\int_X
            \partial \beta = 0\), which implies that \(\exists f \in
            C^{\infty}{(X)}\) such that \(\partial \bar{\partial} f = \partial
            \beta\), hence \(\partial {(\bar{\partial}f - \beta)} = 0\) and thus
            \(\bar{\partial} {(\bar{\beta} - \partial \bar{f})} = 0\). Note that
            the argument \(\bar{\beta} - \partial \bar{f} \in
            H^{{(1,0)}}{(X)}\). We conclude
            \[
              \|\bar{\beta} - \partial \bar{f}\|^2 =  \int_X {\left( \bar{\beta}
                      - \partial \bar{f} \right)} \wedge {\left( \beta -
                      \bar{\partial} f \right)} = \int_X {\left( \bar{\beta} -
              \partial \bar{f} \right)} \wedge \beta = 0
            \]
            hence \(\bar{\beta} - \partial \bar{f} = 0\) which implies that
            \(\bar{\partial}f = \beta\) 
    \end{itemize}
\end{proof}

\subsection{Application: Existence of meromorphic functions}

\begin{theorem}{}
    Let \(X\) be a compact connected Riemann Surface. Let \(p \in X\). Then
    there exists \(f\) meromorphic function on \(X\) non-constant, such that
    \(f\) has a single pole at \(p\) with \(\mathrm{mult}_p {(f)} \le g + 1\) 
\end{theorem}
\begin{proof}{}
    Let \({(U, \varphi )} \ni p\). Let \(q_{0} \in C^{\infty}{(X)}\) such
    that \(\mathrm{Supp}{(q_{0})} \subseteq U \) and \(q_{0} \equiv
    1\) on \(p \in \overline{V} \subseteq U \) and \(V\) open.

    For all \(\mathbf{a}  = {(a_{1}, \dots, a_{g+1} )} \in \mathbb{C}^{g+1}\),
    set \[g_\mathbf{a}{(z)} := q_{0}{(z)} {\left( \sum_{k=1}^{g+1}
    \frac{a_k}{\left[ \varphi {(z)} \right] ^{k}}  \right)} \]
    and extend it by 0 outside \(U\).

    Then \(g_\mathbf{a}\) is meromoprhic on \(V\), but not on \(X\). However, 
    \[
      \bar{\partial} g_\mathbf{a} = {\left( \bar{\partial} q_{0} \right)}
      \sum_{k=1}^{g+q} \frac{a_k}{[\varphi {(z)}]^{k}}
    \]
    and this is all equal to 0 on \(V\), since \(g_\mathbf{a} \equiv 1\). This
    means that \(\bar{\partial} g_\mathbf{a} \in \Omega^{1}{(X)}\).

    The map
    \begin{align*}
        \mathbb{C}^{g+1} &\longrightarrow H^{{(0,1)}}{(X)} \\
         \mathbf{a}  &\longmapsto \left[ \bar{\partial} g_\mathbf{a}  \right] 
    \end{align*}
    is a linear map, and \(\dim H^{{(0,1)}}{(X)} = g < g + 1\), hence there
    exists \(\mathbf{a}  \in \mathbb{C}^{g+1} \sminus \{0\} \) such that
    \(\left[ \bar{\partial} g_\mathbf{a}  \right] = 0\), hence there exists \(f
    \in C^{\infty}{(X)}\) such that \(\bar{\partial} g_\mathbf{a} =
    \bar{\partial} f\) 

    Now observe that \(\mathfrak{f} := g_\mathbf{a} - f\) is holomorphic on \(X \sminus
    \{p\} \) (\(\bar{\partial} \mathfrak{f} = 0\)). Near \(p\) (on \(V\)), this
    can be expressed as \(\underbrace{\mathfrak{f}{(z)}}_{\bar{\partial} = 0}  =
    \underbrace{\sum_{k=1}^{g+1} \frac{a_k}{\left[
    \varphi {(z)}\right] ^{k}}}_{\bar{\partial} = 0}  +
    \underbrace{f{(z)}}_\text{holo near \(p\)}\). Hence, \(\mathfrak{f}\) is
    meromorphic.

\end{proof}

\begin{corollary}{}
    If \(X\) is a compact connected Riemann Surface and \(g{(X)} = 0\), then \(X
    \cong \widehat{\mathbb{C}}\) as Riemann Surfaces.
\end{corollary}
\begin{proof}{}
    Since \(g=0\), there exists a meromorphic map with a single simple pole \(f
    : X \to \widehat{\mathbb{C}}\). Hence \(\deg f = 1\), then \(f\) is an
    isomorphism.
\end{proof}

\begin{eser}{Variant}
    Let \(X\) compact connected Riemann Surface with genus \(g\).
    Show that given \(p_{1}, \dots, p_{g+1} \) different points on \(X\), then
    there exists a meromorphic function such that the set of poles is a subset
    of \(\{p_{1}, \dots, p_{g+1}\} \) and all poles are simple.
\end{eser}



